\section{Exact Euler optimization on a minimum-span face}
\label{sec:cocycle-euler}

The bounded root searches of Section~\ref{sec:cocycle-face} can find surfaces
with different topology at the same minimum piece count, but do not optimize
topology over the whole face. The height cell formula in
Section~\ref{sec:cocycle-reuse} leads to a stronger algorithm: when each global
edge has at least two face incidences, maximizing Euler characteristic on a
certified minimum-span face is an exact, integral, polynomial-size network
optimization. The implementation returns the surface in the existing span
certificate format and a separate arithmetic primal--dual certificate of its
Euler optimum. This solves a restricted optimization problem; the minimum-span
face need not contain a disc or annulus even when the knot is trivial.

\subsection{From cell counts to a nonnegative edge objective}

Orient each global edge $e=(a,b)$ and write its coherent cocycle value after an
integral vertex potential as $x_e(p)=c_e+p_b-p_a$. Let $d_e$ count its incidences with
global triangular faces, with multiplicity: an interior face is included
once, a boundary face once, and each of its three local edges contributes an
incidence even if some are globally identified. For three real numbers, the
span is half the sum of their three absolute pairwise differences. Therefore
the normal cell counts give
\[
  2\chi(S_p)=2D(p)-\sum_e(d_e-2)|x_e(p)|,
  \qquad D(p)=\sum_t\operatorname{span}(h_t+p).
\]
This identity also explains why minimizing the piece count alone does not
maximize Euler characteristic: the edge contribution can vary among tied
minima. It follows from counting cells of the actual coherent level surface,
not from dropping normal matching or quadrilateral conditions.

Fix a previously certified minimum $D_*$. On its entire optimal face,
maximizing $\chi$ is equivalent to minimizing
\[
  F(p)=\sum_e w_e|c_e+p_b-p_a|,\qquad w_e=d_e-2.
\]
The nonnegativity hypothesis $d_e\geq2$ matters. It is checked for every edge;
unsupported negative weights cause rejection, rather than silently changing
the objective. No universal lower-degree claim about arbitrary finite
triangulations is used. The sum $W=\sum_e w_e$ is at most $12N$ for $N$
tetrahedra, because there are at most $4N$ global faces and three incidences
per face. Thus weights here have total magnitude linear in the triangulation
size, independently of the cocycle heights.

The span matching supplies the exact face constraints. If its selected low
and high corners in tetrahedron $t$ are $\ell_t,u_t$, impose, for every corner
$i$,
\[
 p_{v_{\ell_t}}-p_{v_i}\leq h_{t,i}-h_{t,\ell_t},\qquad
 p_{v_i}-p_{v_{u_t}}\leq h_{t,u_t}-h_{t,i}.
\]
These are the same $8N$ difference inequalities used for bounded root
extrema. Writing $z_{t,i}=h_{t,i}+p_{v_i}$, the matching dual gives
\[
 D(p)-D_*=\sum_t\bigl[(\max_i z_{t,i}-z_{t,u_t})
                         +(z_{t,\ell_t}-\min_i z_{t,i})\bigr].
\]
The two extremal gaps per tetrahedron are nonnegative and vanish exactly
when all its displayed inequalities hold. Consequently these inequalities
describe all real and integral potentials attaining $D_*$,
including ties and changes in local normal-disc types. There is no additional
restriction to the root extrema already sampled by the earlier algorithm.

\subsection{An independent lower bound and an exact integral solver}

For each objective edge choose a signed flow $f_e\in[-w_e,w_e]$. For each
constraint $p_v-p_u\leq b_{uv}$ choose a nonnegative flow $\lambda_{uv}$ on
$u\to v$. Require the combined flow to balance at every vertex. Then
\[
\begin{split}
 F(p)&\geq\sum_e f_e(c_e+p_b-p_a)\\
     &=\sum_e f_ec_e-\sum_{uv}\lambda_{uv}(p_v-p_u)\\
     &\geq\sum_e f_ec_e-\sum_{uv}\lambda_{uv}b_{uv}.
\end{split}
\]
The checker validates the primal difference inequalities, both kinds of
flow bounds, vertex balance and equality of this lower bound with $F(p)$.
Equality certifies a global optimum over real potentials, hence also over
integer potentials. This proof needs no trust in the flow solver or in a
floating-point optimization result. The arithmetic checker makes no claim
about triangulation provenance, connectedness or knottedness.

The solver constructs an integral equality witness using a minimum-cost
circulation dual. An objective edge has cost $-c_e$, signed capacity interval
$[-w_e,w_e]$, and stored flow $g_e=f_e+w_e\in[0,2w_e]$. A constraint arc has
cost $b_{uv}$ and capacity $W+1$. Start from the supplied feasible span
potential $p^0$. Set $f_e=w_e$ where $x_e(p^0)>0$, $f_e=-w_e$ where it is
negative, and $f_e=0$ where it vanishes; start all constraint flows at zero.
Every usable residual arc then has nonnegative reduced cost with vertex
potential $p^0$.

This initial flow need not balance. Its total positive imbalance is at most
$W$. A super-source supplies vertices with excess incoming flow and a
super-sink receives the deficits. Their initial potentials are respectively
the maximum and minimum $p^0$ on the relevant vertices, preserving
nonnegative residual costs. Reduced-cost Dijkstra searches and capped
potential updates route the imbalance. Each augmentation sends the full
integral bottleneck, at least one unit. Feasibility is guaranteed by the
balanced signed zero flow: undoing the initial signed edge flow alone gives
a feasible repair within the original capacity intervals.

At most $W$ units are sent in total, so no constraint arc of capacity $W+1$
ever saturates. Its surviving forward residual inequality is exactly the
required primal difference constraint for the final potential. For $w_e>0$, the residual forward and reverse inequalities imply the
appropriate sign of $x_e$: an interior signed flow forces $x_e=0$, while
an endpoint flow can have the corresponding nonzero sign. A positive
constraint flow also leaves its reverse residual arc and forces the
constraint to be tight. These complementary equalities produce equality in
the checker's lower bound. All capacities, distances and potentials remain
integral; no separate rounding step or integrality assumption is needed.

\subsection{Cost, binary encoding and verification boundaries}

There are $O(N)$ nodes and arcs and at most $W\leq12N$ shortest-path
augmentations. The indexed graph computation uses $O(N^2\log(N+1))$
arithmetic and comparison operations and $O(N)$ integer storage entries.
Neither the number of height levels nor the number of normal pieces is
expanded. The private generic weighted-difference solver has an augmentation
bound in the \emph{magnitude} of its weight sum; the polynomial assertion here
uses the proved linear bound on the geometric weights, not a claim of strong
polynomiality for arbitrary binary capacities.

A conservative bit bound is sufficient. Let $B$ bound the input cost and
initial-potential bit lengths, $m=O(N)$ be the number of network nodes,
$C$ bound original absolute arc costs, and $P$ the largest current absolute
potential. A shortest simple residual path has reduced length at most
$(m-1)(C+2P)$. The capped update therefore gives
$P'\leq P+(m-1)(C+2P)$. Iterating at most $12N$ times bounds all arithmetic
lengths by $B+O(N\log(N+1))$ bits. Flows have $O(\log(N+1))$ bits. Together
with polynomial-time integer arithmetic, geometry validation and conservative
Python-container accounting, this establishes a bit-polynomial bound for
the complete restricted stage. The tighter indexed operation count is not
presented as a literal Python bit-operation count.

The public optimizer validates the triangulation, signed global coherence
and supplied span certificate before constructing this model. It checks its
new primal--dual equality, replays the returned span certificate, and compares
the optimum edge objective with an independent normal-coordinate Euler count.
The surface still uses \code{normal-cocycle-span-v1}. An unknot claim must
pass the existing independent source-bound disc or annulus verifier; those
verifiers do not import the new optimizer. A negative Euler optimum is a
certificate about this face alone, not a knot certificate. Even maximizing
Euler here does not maximize the number or essentiality of boundary circles
needed by the separate planar-capping search.

\subsection{Source controls and a certified limitation}

The audit uses the preceding 84-source corpus plus the three-crossing unknot
control, with a shared allowance of two million guards per source covering
exterior construction, cocycle extraction, span optimization, Euler
optimization and independent replays. It completes 83 of 85 cases. Gordian
exhausts the allowance during span optimization; \code{circle-33} exhausts it
during Euler optimization. These are incomplete arithmetic experiments, not
negative knot verdicts. All completed geometric models have strictly positive
edge weights; the smallest observed face incidence count is three.

Ten labels improve their initial Euler value: survivor cases 01 and 03,
mirror case 03, Kinoshita--Terasaka, and random cases 0, 19, 30, 34, 142 and 158.
Most gain one Euler unit; random case 158 improves from $-9$ to $-7$.
Random case 142 improves from an annulus to a disc, preserving its certified
122-piece minimum. On all 74 sources covered by the preceding face survey,
the exact optimum agrees with the best Euler value among the 787 previously
recorded root candidates, including each original optimum. Agreement on this
corpus is an observation; the new global dual is what proves optimality over
the whole face.

The 18 nonnegative optima produce 13 distinct source-bound positive
certificates, all independently replayed. No additional source becomes
recognized relative to the current four-tree native policy with annulus
capping. In fact, five already recognized sources---random cases 29, 79, 99,
124 and 150---have a negative optimum on their minimum-span face. Their
existing earlier witnesses lie outside that face. The exact optimizer is
therefore not a replacement for the broader existing candidate schedule.
Regina independently confirms connectedness, orientability, Euler
characteristic and the compressing-disc predicate on the 75 sources with at
most twelve crossings, including the extra unknot control. The audit retains
all 73 distinct indexed arithmetic models and duals for direct replay without
running the new optimizer.

For the closure of the two-strand braid $[1,1,-1]$, cancellation leaves the
one-crossing unknot. Nevertheless, the span minimum is 43 normal pieces and
the new dual proves that every surface in that face has Euler characteristic
at most $-1$. The trefoil control has the same Euler upper bound, at a
45-piece minimum. This gives a concrete obstruction to treating exhaustive
minimum-span-face optimization as complete disc or annulus discovery. It
also separates this result from a minimum-genus or Thurston-norm algorithm:
the fixed piece-count optimum is an additional restriction, and the separate
planar-capping condition is not optimized by maximizing Euler characteristic.

The arithmetic tests compare 240 small random constrained problems against
exhaustive integer grids, including loops, parallel edges, zero weights and
repeated constraints. Exact primal--dual equality additionally certifies the
real optimum in each case. Scaling costs and bounds by $2^{20000}$ preserves
all flow-operation statistics while scaling the objective, and the resulting
certificate survives binary JSON transport. Further checks cover independent
replay with the producer disabled, malformed duals, incoherent heights,
negative geometric weights, hexadecimal certificate fields, cancellation,
source-bound positive verification and exact work-allowance boundaries.

\subsection{Validation, timings and integration scope}

All 14 focused tests pass in 1.753 seconds; the full maintained suite passes
all 1,219 tests in 254.202 seconds. The final source audit takes 19.307
seconds. The audit and full suite partly overlap, so those elapsed times
are correctness-run observations rather than performance comparisons. The
final audit and both timing runs pin the same 501 source files. Existing
source-bound disc, annulus, planar, span and weighted-orbit verifiers remain
byte-identical to the baseline revision \code{3c0817c50}.

The performance endpoint starts from a diagram and includes finite exterior
construction, primitive cocycle extraction, minimum-span optimization, the
selected search and independent source-bound replay of its selected surface.
The old arm searches all distinct root extrema, using the number of global
vertices as its root allowance, and retains the highest Euler value found.
The new arm solves the whole face and checks its arithmetic dual. This is a
complete \emph{candidate-construction} pipeline, not complete knot
recognition. Earlier native successes are deliberately not allowed to skip
the experiment, and no exact knot-recognition fallback is timed here.

The first five-round run contains 120 measured calls and 24 warm-ups and
takes 24.281 seconds. Variation in the original-disc and Kinoshita--Terasaka
cases motivates the retained nine-round repeat: 216 measured calls and 24
warm-ups in 40.195 seconds. Each round randomizes roots, roots A/A, exact and
exact A/A calls. No tests, audits, builds or commits overlap either timing
run. Every measured call completes, and all per-arm Euler values, piece
counts, operation counts and certificate hashes are stable within and
between both runs. Primary-arm median times and median paired ratios from
the final run follow; paired ratios need not equal quotients of the displayed
medians.

\begin{center}
\small
\begin{tabular}{lrrrrr}
\toprule
Input & Roots (ms) & Exact (ms) & Roots/exact & Roots A/A & Exact A/A \\
\midrule
Trefoil & 38.209 & 37.957 & 1.021 & 0.940 & 0.976 \\
Original disc & 48.704 & 71.274 & 0.690 & 1.004 & 1.003 \\
Tied disc & 115.241 & 115.147 & 0.989 & 0.986 & 0.990 \\
Random 19 & 228.711 & 218.924 & 1.063 & 1.001 & 1.002 \\
Random 158 & 333.819 & 259.215 & 1.283 & 0.996 & 0.994 \\
Kinoshita--Terasaka & 202.822 & 227.087 & 0.825 & 1.027 & 0.995 \\
\bottomrule
\end{tabular}

\end{center}

Random case 158 improves by paired factors 1.272 and 1.283 in the two runs;
random case 19 improves by 1.054 and 1.062. The original-disc case is slower,
with ratios 0.677 and 0.690, and Kinoshita--Terasaka is slower with ratios
0.835 and 0.825. Those final ratios correspond to about 45 and 21 percent
additional time. The tied-disc case is near parity. The trefoil's small
apparent improvement is not evidence of a reliable gain: its final roots
A/A ratio is 0.940. Absolute Kinoshita--Terasaka times also vary substantially
between runs. The data support mixed costs, not a uniform speedup.

The new callable \code{maximize\_cocycle\_face\_euler} API is consequently
available for explicit candidate construction and research; it is not added
to the default native schedule. All existing candidate and portfolio defaults
remain unchanged. Its theoretical gain is an exact optimum over a face described by $8N$
inequalities, together with concrete certified limitations of that face. It supplies neither all normal-surface candidates
nor a complete planar-capping search, a hierarchy algorithm, or a general
subexponential or quasipolynomial recognition bound.

Reproduce with \code{normal\_orbit\_research.euler} in \code{audit} mode or
\code{benchmark --rounds 9} mode. The indexed models and arithmetic duals
are retained in the audit JSON for direct independent replay. The incoming
unknot archive inventory and the checked \code{openai/math} head
\code{fd4aeeb2e} remain unchanged. The new reduction is derived from the
maintained height-cell identity and matching-face description above.


